% Provenance (never printed): the public-finance strand of ../economics_series.md (grade 8: what governments spend
% on; the main taxes; budget, deficit and debt; a tax changes behaviour, first look; government failure named once,
% taught in g11). Brief: sources/book1/BRIEFS.md, g8-07. Chart data: figdata/grade-8/07-taxes-and-public-spending
% (JST R6 expenditure, gdp, debtgdp; Brenland TAX_BANDS and accounts).
% Ledger: g8:cigarette-tax-elasticity, g8:vat-rate-2026.
\chapter{Taxes and Public Spending}\label{ch:g8:taxes-and-public-spending}

\begin{openingchart}{Central government spending as a share of GDP in six countries, 1880--2020, in
percent. Regional and local governments are not included.}{Jord\`a--Schularick--Taylor
Macrohistory Database, central government expenditure and GDP}
\begin{tikzpicture}
\begin{axis}[omtime, xmin=1880, xmax=2020, ymin=0, ymax=70, xlabel={Year},
  ylabel={Percent of GDP}, xtick={1880,1900,1920,1940,1960,1980,2000,2020},
  legend style={at={(0.5,-0.2)}, anchor=north, legend columns=6}]
\addplot table[x=year, y=usa] {figdata/grade-8/07-taxes-and-public-spending/out/public_spending.csv};
\addplot table[x=year, y=gbr] {figdata/grade-8/07-taxes-and-public-spending/out/public_spending.csv};
\addplot table[x=year, y=fra] {figdata/grade-8/07-taxes-and-public-spending/out/public_spending.csv};
\addplot table[x=year, y=swe] {figdata/grade-8/07-taxes-and-public-spending/out/public_spending.csv};
\addplot table[x=year, y=jpn, unbounded coords=jump] {figdata/grade-8/07-taxes-and-public-spending/out/public_spending.csv};
\addplot table[x=year, y=deu, unbounded coords=jump] {figdata/grade-8/07-taxes-and-public-spending/out/public_spending.csv};
\legend{United States, United Kingdom, France, Sweden, Japan, Germany}
\end{axis}
\end{tikzpicture}
\end{openingchart}

% figdata: facts (spend_usa_1913 1.8, 1950 14.2, 2019 20.8; spend_gbr_1913 8.0, 1950 26.5, 2019 36.8;
% spend_fra_1918 63.0; spend_gbr_1944 62.8; spend_usa_1944 40.7)
In 1913, the central government of the United States spent the equivalent of $1.8\,\%$ of GDP; in
2019, $20.8\,\%$. In the United Kingdom, the share went from $8.0\,\%$ to $36.8\,\%$. Twice it jumped:
during the two world wars, France's central government spent $63.0\,\%$ of GDP in 1918, the United
Kingdom's $62.8\,\%$ in 1944. After each war the share fell back, but not to where it had been.

What do governments spend on, where does the \omterm{def:g7:money:money}{money} come from, and what happens when they spend more
than they collect?

\begin{recall}
\omterm{def:g7:households:household}{Households} receive \omterm{def:g7:households:transfer}{transfers} and pay \omterm{def:g8:taxes-and-public-spending:tax}{taxes}; \omterm{def:g7:households:disposable-income}{disposable income} is gross \omterm{def:g7:households:income}{income} minus \omterm{def:g8:taxes-and-public-spending:tax}{taxes} plus
\omterm{def:g7:households:transfer}{transfers} (\cref{ch:g7:households}). GDP is the value of everything produced in a country in a year
(\cref{ch:g8:gdp}). A borrower pays \omterm{def:g8:banks-saving-interest:interest}{interest} on a \omterm{def:g8:banks-saving-interest:loan}{loan} (\cref{ch:g8:banks-saving-interest}).
\end{recall}

\section{What governments spend on}

\begin{definition}[Public spending]\label{def:g8:taxes-and-public-spending:public-spending}
\emph{Public spending}\index{public spending} is the spending of governments: on the goods and
services they provide, on the \omterm{def:g7:households:wage}{wages} of the people they employ, on \omterm{def:g7:households:transfer}{transfers} to \omterm{def:g7:households:household}{households} and \omterm{def:g7:firms:firm}{firms},
and on \omterm{def:g8:banks-saving-interest:interest}{interest} on their debt.
\end{definition}

Governments run schools, hospitals, courts, the police and the army; they build and repair roads;
they pay pensions and support people without work. In many rich countries, pensions, health and
education are among the largest items. Part of the spending buys goods and services that the government
provides; another part only passes \omterm{def:g7:money:money}{money} from some \omterm{def:g7:households:household}{households} to others, as \omterm{def:g7:households:transfer}{transfers}.

\begin{omfigure}
\includegraphics[width=0.6\linewidth]{images/book1/ai/g8-taxes-town-street.jpg}

{\small A school, a road being repaired and a health clinic: three things paid for, in many countries,
by \omterm{def:g8:taxes-and-public-spending:public-spending}{public spending}.}
\end{omfigure}

\begin{measure}[General and central government]
The \emph{central government} is the national government. The \emph{general government} adds the
regions, the towns and the social-security funds that pay pensions and health care. In a country
where regions run the schools and the hospitals, central government spending is a small part of the
total: Germany's line in the opening chart is low for that reason, not because German governments
spend little. Comparing countries therefore needs general government figures; in the database of the
opening chart, the long series exist only for the central government.
\end{measure}

\section{The main taxes}

\begin{definition}[Tax]\label{def:g8:taxes-and-public-spending:tax}
A \emph{tax}\index{tax} is a payment that people or \omterm{def:g7:firms:firm}{firms} must make to a government by law, without
receiving a particular good or service in exchange.
\end{definition}

\begin{definition}[Income tax]\label{def:g8:taxes-and-public-spending:income-tax}
An \emph{income tax}\index{income tax} is a \omterm{def:g8:taxes-and-public-spending:tax}{tax} computed on the \omterm{def:g7:households:income}{income} of a person or a \omterm{def:g7:households:household}{household}.
\end{definition}

\begin{definition}[Tax on purchases]\label{def:g8:taxes-and-public-spending:tax-on-purchases}
A \emph{tax on purchases}\index{tax on purchases} is a \omterm{def:g8:taxes-and-public-spending:tax}{tax} added to the \omterm{def:g7:money:price}{price} of goods and services
when they are bought. A \emph{sales \omterm{def:g8:taxes-and-public-spending:tax}{tax}} is charged once, when the final buyer pays; a
\emph{value-added \omterm{def:g8:taxes-and-public-spending:tax}{tax}} (VAT) is charged at each stage of production on the \omterm{def:g8:gdp:value-added}{value added} at that stage.
\end{definition}

\begin{definition}[Payroll tax]\label{def:g8:taxes-and-public-spending:payroll-tax}
A \emph{payroll tax}\index{payroll tax} is a \omterm{def:g8:taxes-and-public-spending:tax}{tax} computed on the \omterm{def:g7:households:wage}{wages} that employers pay, often used to
pay for pensions and health care.
\end{definition}

\begin{dated}{2026}
% ledger: g8:vat-rate-2026
In the United Kingdom, the standard rate of VAT is $20\,\%$ (since 4 January 2011): a good priced at
$100$ pounds before \omterm{def:g8:taxes-and-public-spending:tax}{tax} costs $120$ pounds in the shop.
\end{dated}

Most \omterm{def:g8:taxes-and-public-spending:income-tax}{income taxes} are cut into bands, with a higher rate on each higher band. Brenland's \omterm{def:g8:taxes-and-public-spending:income-tax}{income tax}
of 2025 is computed on a \omterm{def:g7:households:household}{household}'s gross \omterm{def:g7:households:income}{income}: nothing on the first $12\,000$ dalens, $20\,\%$ on
the part between $12\,000$ and $40\,000$, $35\,\%$ on the part between $40\,000$ and $100\,000$, and
$45\,\%$ on the part above $100\,000$.

\begin{omchart}{Brenland's income tax, 2025: the rate paid on each additional dalen of gross household
income.}{Brenland (invented)}
\begin{tikzpicture}
\begin{axis}[omrecord, xmin=0, xmax=140000, ymin=0, ymax=50, height=0.4\linewidth,
  xlabel={Gross household income (dalens a year)}, ylabel={Rate (percent)},
  xtick={0,12000,40000,100000,140000}, scaled x ticks=false,
  xticklabels={0, 12\,000, 40\,000, 100\,000, 140\,000}]
\addplot[omDef, very thick, const plot] table[x=income, y=rate] {figdata/grade-8/07-taxes-and-public-spending/out/tax_bands.csv};
\end{axis}
\end{tikzpicture}
\end{omchart}

\begin{definition}[Average tax rate]\label{def:g8:taxes-and-public-spending:average-tax-rate}
The \emph{average tax rate}\index{average tax rate} of a \omterm{def:g7:households:household}{household} is the \omterm{def:g8:taxes-and-public-spending:tax}{tax} it pays divided by its
\omterm{def:g7:households:income}{income}, in percent.
\end{definition}

\begin{definition}[Marginal tax rate]\label{def:g8:taxes-and-public-spending:marginal-tax-rate}
The \emph{marginal tax rate}\index{marginal tax rate} of a \omterm{def:g7:households:household}{household} is the rate of \omterm{def:g8:taxes-and-public-spending:tax}{tax} on one more
dalen of \omterm{def:g7:households:income}{income}: the rate of the band its \omterm{def:g7:households:income}{income} is in.
\end{definition}

% figdata: facts (tax_30000 3600, tax_60000 12600)
A Brenland \omterm{def:g7:households:household}{household} with a gross \omterm{def:g7:households:income}{income} of $30\,000$ dalens pays $(30\,000 - 12\,000) \times 0.20 =
3\,600$ dalens. Its \omterm{def:g8:taxes-and-public-spending:average-tax-rate}{average tax rate} is $3\,600 \div 30\,000 = 12\,\%$, but its marginal rate is
$20\,\%$: of each extra $100$ dalens it earns, $20$ go in \omterm{def:g8:taxes-and-public-spending:tax}{tax}. The average rate is below the marginal rate
here for every \omterm{def:g7:households:income}{income} above $12\,000$ dalens, because the first $12\,000$ dalens are not taxed.

\section{The budget, the deficit and the debt}

\begin{definition}[Government budget]\label{def:g8:taxes-and-public-spending:government-budget}
A \emph{government budget}\index{government budget} is the list of a government's \omterm{def:g7:firms:revenue}{revenue} (mostly
\omterm{def:g8:taxes-and-public-spending:tax}{taxes}) and spending for a year.
\end{definition}

\begin{definition}[Budget deficit]\label{def:g8:taxes-and-public-spending:budget-deficit}
A \emph{budget deficit}\index{budget deficit} is the amount by which spending is larger than \omterm{def:g7:firms:revenue}{revenue}
in a year.
\end{definition}

\begin{definition}[Budget surplus]\label{def:g8:taxes-and-public-spending:budget-surplus}
A \emph{budget surplus}\index{budget surplus} is the amount by which \omterm{def:g7:firms:revenue}{revenue} is larger than spending
in a year.
\end{definition}

\begin{definition}[Public debt]\label{def:g8:taxes-and-public-spending:public-debt}
\emph{Public debt}\index{public debt} is the total that a government owes to its lenders at a given
date: the sum of its past deficits minus its past surpluses.
\end{definition}

A government with a deficit borrows the difference, usually by selling bonds
(\cref{ch:g8:banks-saving-interest}). The debt is like the water in a bathtub: deficits pour in,
surpluses drain it, and each year's debt is last year's debt plus this year's deficit.

\begin{omfigure}
\begin{tikzpicture}[font=\small]
  \draw[very thick] (0,3) -- (0,0) -- (5,0) -- (5,3);
  \fill[omDef!35] (0.05,0.05) rectangle (4.95,1.8);
  \node at (2.5,0.9) {public debt};
  \draw[thick, -{Stealth[length=3mm]}] (1.2,4.2) -- (1.2,2.1);
  \node[align=center, anchor=south] at (1.2,4.2) {deficit\\(spending $>$ revenue)};
  \draw[thick] (4.4,0) -- (4.4,-0.5) -- (6.2,-0.5);
  \draw[thick, -{Stealth[length=3mm]}] (6.2,-0.5) -- (7.0,-0.5);
  \node[align=left, anchor=west] at (7.1,-0.5) {surplus\\(revenue $>$ spending)};
\end{tikzpicture}

{\small The debt as a bathtub: a deficit adds to the debt, a surplus takes away from it.}
\end{omfigure}

Debt is usually measured as a share of GDP, because a country with a larger GDP can carry a larger
debt: the same debt is a smaller burden when the economy and the \omterm{def:g8:taxes-and-public-spending:tax}{tax} \omterm{def:g7:firms:revenue}{revenue} it brings grow.

% figdata: facts (debt_gbr_max 269.8 1946; debt_gbr_1919 142.8; debt_usa_1946 118.9; debt_usa_2020 128.1;
% debt_jpn_2020 253.9; debt_fra_max 237.0 1921)
\begin{omchart}{Public debt as a share of GDP in six countries, 1880--2020, in
percent.}{Jord\`a--Schularick--Taylor Macrohistory Database, \omterm{def:g8:taxes-and-public-spending:public-debt}{public debt} to GDP ratio}
\begin{tikzpicture}
\begin{axis}[omtime, xmin=1880, xmax=2020, ymin=0, ymax=280, xlabel={Year},
  ylabel={Percent of GDP}, xtick={1880,1900,1920,1940,1960,1980,2000,2020},
  legend style={at={(0.5,-0.2)}, anchor=north, legend columns=6}]
\addplot table[x=year, y=usa] {figdata/grade-8/07-taxes-and-public-spending/out/debt_gdp.csv};
\addplot table[x=year, y=gbr] {figdata/grade-8/07-taxes-and-public-spending/out/debt_gdp.csv};
\addplot table[x=year, y=fra, unbounded coords=jump] {figdata/grade-8/07-taxes-and-public-spending/out/debt_gdp.csv};
\addplot table[x=year, y=swe] {figdata/grade-8/07-taxes-and-public-spending/out/debt_gdp.csv};
\addplot table[x=year, y=jpn, unbounded coords=jump] {figdata/grade-8/07-taxes-and-public-spending/out/debt_gdp.csv};
\addplot table[x=year, y=deu, unbounded coords=jump] {figdata/grade-8/07-taxes-and-public-spending/out/debt_gdp.csv};
\legend{United States, United Kingdom, France, Sweden, Japan, Germany}
\end{axis}
\end{tikzpicture}
\end{omchart}

Wars built the largest debts: the United Kingdom's reached $269.8\,\%$ of GDP in 1946, the United
States' $118.9\,\%$. In the following decades both shares fell a long way, as GDP grew faster than
the debt. In 2020, the share was $128.1\,\%$ in the United States and $253.9\,\%$ in Japan, without a
war.

\section{A tax changes behaviour}

\begin{model}{Taxes change behaviour}
\assumes People respond to \omterm{def:g7:money:price}{prices} and \omterm{def:g7:households:wage}{wages} after \omterm{def:g8:taxes-and-public-spending:tax}{tax}.
\predicts A \omterm{def:g8:taxes-and-public-spending:tax}{tax} on a good raises its \omterm{def:g7:money:price}{price} and lowers the quantity bought, more so when buyers have
other goods they can switch to.
\tested Studies of cigarette \omterm{def:g8:taxes-and-public-spending:tax}{taxes} and \omterm{def:g7:money:price}{prices} in rich countries find that a $10\,\%$ rise in the \omterm{def:g7:money:price}{price}
of cigarettes lowers the quantity smoked by about $4\,\%$; most estimates lie between 2 and $6\,\%$
(National Cancer Institute and World Health Organization, 2016).
\breaks In the short run, when buyers have no substitute or are addicted, the quantity hardly moves;
when the \omterm{def:g8:taxes-and-public-spending:tax}{tax} can be evaded, recorded sales fall while buying from untaxed sellers rises.
\end{model}
% ledger: g8:cigarette-tax-elasticity

A \omterm{def:g8:taxes-and-public-spending:tax}{tax} raises \omterm{def:g7:money:money}{money}, and it changes what people buy, sell and do. Both effects matter: a \omterm{def:g8:taxes-and-public-spending:tax}{tax} that
makes people buy much less of a good also raises less \omterm{def:g7:money:money}{money} than a simple sum would suggest. How much
the quantity falls is measured with the elasticity of \cref{ch:g9:elasticity-first-look}.

\section{Who decides}

\omterm{def:g8:taxes-and-public-spending:tax}{Taxes} and spending are decided by parliaments and governments, through elections and laws. Those who
decide face their own pressures, from voters, from groups that gain from a \omterm{def:g8:taxes-and-public-spending:tax}{tax} or a spending item, and
from the next election, and they can get things wrong just as \omterm{def:g7:markets-and-prices:market}{markets} can. How public decisions are
made, and when they fail, is studied in the volume for grades 10 to 12.

\section{Exercises}

\begin{exercise}[$\star$]\label{exo:g8:taxes-and-public-spending:1}
Use the text. What share of GDP did central government spend in 1913 and in 2019 in the United States?
In the United Kingdom?
\end{exercise}

\begin{exercise}[$\star$]\label{exo:g8:taxes-and-public-spending:2}
A government collects $450$ million dalens and spends $500$ million. Compute its deficit.
\end{exercise}

\begin{exercise}[$\star$]\label{exo:g8:taxes-and-public-spending:3}
A government's debt is $1\,200$ million dalens. This year its deficit is $50$ million. What is its debt
at the end of the year?
\end{exercise}

\begin{exercise}[$\star$]\label{exo:g8:taxes-and-public-spending:4}
A sales \omterm{def:g8:taxes-and-public-spending:tax}{tax} of $20\,\%$ is added to a \omterm{def:g7:money:price}{price} of $15$ dalens. What does the buyer pay?
\end{exercise}

\begin{exercise}[$\star\star$]\label{exo:g8:taxes-and-public-spending:5}
Compute the \omterm{def:g8:taxes-and-public-spending:income-tax}{income tax} of a Brenland \omterm{def:g7:households:household}{household} with a gross \omterm{def:g7:households:income}{income} of $45\,000$ dalens.
\end{exercise}

\begin{exercise}[$\star\star$]\label{exo:g8:taxes-and-public-spending:6}
For the \omterm{def:g7:households:household}{household} of exercise 5, compute the \omterm{def:g8:taxes-and-public-spending:average-tax-rate}{average tax rate} (to one decimal) and give the \omterm{def:g8:taxes-and-public-spending:marginal-tax-rate}{marginal
tax rate}.
\end{exercise}

\begin{exercise}[$\star\star$]\label{exo:g8:taxes-and-public-spending:7}
Use the debt chart. In which years did the United Kingdom's and France's debts peak, and what was
happening then?
\end{exercise}

\begin{exercise}[$\star\star$]\label{exo:g8:taxes-and-public-spending:8}
A \omterm{def:g8:taxes-and-public-spending:tax}{tax} raises the \omterm{def:g7:money:price}{price} of cigarettes by $20\,\%$. Using the estimate of the model box, by about what
percentage does the quantity smoked fall? If $1\,000$ million packs were sold before, about how many
are sold after?
\end{exercise}

\begin{exercise}[$\star\star$]\label{exo:g8:taxes-and-public-spending:9}
After a large \omterm{def:g8:taxes-and-public-spending:tax}{tax} rise on cigarettes, recorded sales fall by $15\,\%$, but surveys show that people
smoke about as much as before. Which part of the model box does this match, and what might explain
it?
\end{exercise}

\begin{exercise}[$\star\star\star$]\label{exo:g8:taxes-and-public-spending:10}
In 2025, Brenland's debt was $236\,818$ million dalens and its GDP $393\,264$ million. Suppose the
deficit is $8\,000$ million in each of the next two years and GDP grows by $4\,\%$ a year. Compute the
debt as a share of GDP in 2025 and after two years. Did the share rise or fall?
\end{exercise}

\begin{exercise}[$\star\star\star$]\label{exo:g8:taxes-and-public-spending:11}
A worker earns $3\,000$ dalens a month. A \omterm{def:g8:taxes-and-public-spending:payroll-tax}{payroll tax} of $10\,\%$ of the \omterm{def:g7:households:wage}{wage}, paid by the employer,
is introduced. After some time, the worker's \omterm{def:g7:households:wage}{wage} is $2\,800$ dalens. (a) How much \omterm{def:g8:taxes-and-public-spending:tax}{tax} does the
employer pay? (b) How much does the worker cost the employer in all, before and after? (c) How much of
the \omterm{def:g8:taxes-and-public-spending:tax}{tax} is borne by the worker, through a lower \omterm{def:g7:households:wage}{wage}, and how much by the employer?
\end{exercise}

\begin{exercise}[$\star\star\star$]\label{exo:g8:taxes-and-public-spending:12}
Explain what the opening chart cannot tell about what the spending achieved.
\end{exercise}

\section{Problem: Brenland's Budget}

\begin{problem}[{Weekend problem --- Brenland's revenue, spending and deficit, its public debt, and a
household's income tax}]
\label{pb:g8:taxes-and-public-spending:1}
% figdata: budget (2015 248438 99375 104344 4969 159977; 2019 289032 115613 121393 5780 181790;
% 2020 279869 109149 128740 19591 201381; 2025 393264 157306 165171 7865 236818)
Everything in this problem is invented. Brenland's public accounts, in millions of dalens:
\begin{center}
\begin{tabular}{lrrrr}
\toprule
Year & 2015 & 2019 & 2020 & 2025 \\
\midrule
GDP & 248\,438 & 289\,032 & 279\,869 & 393\,264 \\
\omterm{def:g7:firms:revenue}{Revenue} & 99\,375 & 115\,613 & 109\,149 & 157\,306 \\
Spending & 104\,344 & 121\,393 & 128\,740 & 165\,171 \\
Debt at the end of the year & 159\,977 & 181\,790 & 201\,381 & 236\,818 \\
\bottomrule
\end{tabular}
\end{center}

\medskip
\noindent\textbf{Part I --- The deficit.}
\begin{enumerate}
  \item Compute the deficit in 2015.
  \item Compute the deficit in 2020. Using the table, give two reasons why it was so much larger than in
        2019.
  \item Compute the deficit in 2025, and as a share of GDP.
  \item Compute \omterm{def:g8:taxes-and-public-spending:public-spending}{public spending} as a share of GDP in 2025.
\end{enumerate}

\medskip
\noindent\textbf{Part II --- The debt.}
\begin{enumerate}[resume]
  \item Check that the debt of 2020 equals the debt of 2019 plus the deficit of 2020.
  \item Compute the debt as a share of GDP in 2019 and in 2020.
  \item Why did the share rise so much in 2020?
  \item Compute the debt as a share of GDP in 2025.
  \item If the government pays $3\,\%$ \omterm{def:g8:banks-saving-interest:interest}{interest} on its debt of 2025, how much \omterm{def:g8:banks-saving-interest:interest}{interest} does it pay in a
        year?
\end{enumerate}

\medskip
\noindent\textbf{Part III --- A household's income tax.} Use Brenland's \omterm{def:g8:taxes-and-public-spending:tax}{tax} bands.
\begin{enumerate}[resume]
  \item Compute the \omterm{def:g8:taxes-and-public-spending:tax}{tax} of a \omterm{def:g7:households:household}{household} with a gross \omterm{def:g7:households:income}{income} of $30\,000$ dalens.
  \item Compute the \omterm{def:g8:taxes-and-public-spending:tax}{tax} of a \omterm{def:g7:households:household}{household} with $60\,000$ dalens.
  \item Compute both \omterm{def:g7:households:household}{households}' \omterm{def:g8:taxes-and-public-spending:average-tax-rate}{average tax rates}.
  \item What is the \omterm{def:g8:taxes-and-public-spending:marginal-tax-rate}{marginal tax rate} of the second \omterm{def:g7:households:household}{household}? How much more \omterm{def:g8:taxes-and-public-spending:tax}{tax} does it pay if it
        earns $1\,000$ dalens more?
\end{enumerate}
\end{problem}
